% !TEX root = ../../main-es.tex
% appendix/es/A-criterios.tex
\chapter{Correlación objetivos \ensuremath{\leftrightarrow} criterios EICT}
\label{app:criterios}

\porque{Trazabilidad institucional.}{%
  Tabla auto-generada por cada llamada a \texttt{\textbackslash criterio} a lo largo del
  documento. Si una sección falta aquí, no está etiquetada. La correlación con
  \texttt{docs/institucion/normativa/criterios-tesis.md} vive \emph{dentro} del PDF y no se desincroniza a mano.}

\section*{Mapa de secciones por dimensión EICT}
\listacriterios

\section{Análisis de viabilidad por objetivo}
\porque{Reescritura de docs/institucion/normativa/criterios-tesis.md §"Análisis de Viabilidad".}{%
  Para O1--O5 (no los viejos O1--O5 de atención). Formato: una fila por objetivo con
  métrica, criterio EICT afectado y viabilidad académica. Las fases administrativas y
  la rúbrica de \texttt{docs/institucion/normativa/criterios-tesis.md} no se repiten aquí: son hechos institucionales.}

\begin{longtable}{p{0.10\textwidth} p{0.30\textwidth} p{0.25\textwidth} p{0.30\textwidth}}
\toprule
\textbf{Obj.} & \textbf{Métrica} & \textbf{Criterio EICT} & \textbf{Viabilidad} \\
\midrule
O1 & Semántica total; leyes verificadas & Rigor metodológico & Alta: lectura literal del código existente \\
O2 & Demostración revisada & Fundamentación teórica & Media: requiere revisor con perfil formal \\
O3 & Soundness + 0 FP + cobertura & Plan de acción / rigor & Alta: esquema de salida ya presente en el motor \\
O4 & OR e IC 95\% pre-declarados & Análisis crítico & Media-alta: arnés con ablación \\
O5 & Arnés versionado M1; M2/M3 especificadas & Reproducibilidad & Alta \\
\bottomrule
\end{longtable}
